\section{Agent B: the LLM agent}
\label{sec:llm}

Agent~B is Agent~A plus a thin LLM layer. The LLM is used \emph{only} for
high-level interpretation --- turning a natural-language mission into structured
belief updates --- never inside the tick loop and never to drive movement, which
stays deterministic BFS (Section~\ref{sec:arch}). This keeps reactivity intact
while letting the agent understand requests it was not hard-coded for.

\paragraph{Tools.} A small validated tool registry is defined ---
\texttt{get\_state} (read beliefs), \texttt{set\_mission} (install a structured
mission), \texttt{block\_tiles} / forbidden-delivery setters, and
\texttt{send\_to\_teammate} / \texttt{reply\_to\_agent} (communication) ---
available for tool-loop experiments and future extensions (the lab8 pattern).
The production path does not drive a live tool-calling loop: it interprets each
mission into schema-constrained JSON (next paragraph). The design choice is
deliberate either way --- there is \emph{no} low-level movement tool, so the LLM
expresses \emph{what} the mission is and the BDI machinery decides \emph{how} ---
and every tool argument is schema-validated before it touches beliefs.

\paragraph{Prompt engineering and fallback.} The mission-interpretation prompt
asks for schema-constrained JSON (mission kind, targets, tolerance, thresholds);
the result is validated and rejected if malformed. A deterministic parser backs
the LLM for latency- and availability-critical templates --- JSON coordinate
lists, red/green light, arithmetic question--answer, and value-threshold wording
(``lower or equal to'', ``Threshold is''). The fallback is a \emph{resilience
layer, not the intelligence}: an all-fallback pass (gateway unreachable) still
passed every scenario, but the LLM is what generalizes to reworded requests.

\paragraph{Request handling across Challenge~2 levels.} \emph{Level~1 (atomic
requests)}: arithmetic and positional questions are answered directly
(e.g.\ 26c2\_3, ``compute~$\to$~22'', rewarded). \emph{Level~2 (strategy
adaptation)}: missions install constraints/policies that reshape the utility ---
positional goals (\texttt{go\_to}), \texttt{deliver\_at} a tile,
\texttt{deliver\_exactly\_n}, and \texttt{deliver\_less\_value\_than}. Several
behaviours needed live fixes that the report records as findings:
\texttt{deliver\_at} waits briefly before putdown so the observer can award the
target delivery; \texttt{deliver\_exactly\_n} suppresses premature deliveries
until the required batch is carried, with a final pre-putdown guard for missions
that arrive mid-intention; forbidden-tile updates invalidate stale paths before
the next move; and \texttt{deliver\_less\_value\_than} uses selective under-cap
putdown (deliver only the parcels below the threshold).

\paragraph{Latency safety --- a live finding.} With the LLM in the loop,
interpretation latency is variable and can be large (1.8--8.7\,s observed). On
26c2\_4 (movement through a tile band forbidden under a $-1000$ penalty) an
8.7\,s interpretation let the agent cross a forbidden tile and take the penalty,
scoring $-377$. The fix parses each incoming message deterministically and
\emph{pre-applies} any safety-critical constraint (forbidden \texttt{go\_to} /
\texttt{deliver\_at}, light state; \texttt{MissionInterpreter.isSafetyCritical})
to beliefs the instant it arrives; the LLM result still runs and remains
authoritative, reconciling idempotently. This closes the latency window while
keeping the LLM the primary interpreter: the \texttt{mission\_preapplied} log
precedes the LLM completion by $\sim$3\,s, and the same scenario then scored
$+648$ with no penalty (\texttt{source=llm} preserved). Malformed LLM outputs
are caught by validation and fall back to the deterministic parser. The full
validation is reported in Section~\ref{sec:experiments}.
